\documentclass[10pt,a4paper]{amsart}
\usepackage{amsmath,amssymb,amsthm,graphicx}
% Retained preamble for the Stacks Project chapter snapshot in this directory.
% The chapter's own preamble.tex is a copy of the shared Stacks Project preamble; the block
% below is sliced MECHANICALLY, line for line, from it. Nothing was added, removed or
% redefined; the wrapper class, the xr-hyper/hyperref wiring and the externaldocument list
% are omitted because the assembler supplies the document class.
%
% lines taken from preamble.tex: 3, 5, 9, 11, 12, 156, 157, 158, ... (51 lines total)
%
% package coverage restored from the chapter's own preamble (2026-10-13): xy, xy-2cell

\usepackage[all]{xy}
\xyoption{2cell}
\UseAllTwocells

\newenvironment{reference}{\comment}{\endcomment}
\newenvironment{slogan}{\comment}{\endcomment}
\newenvironment{history}{\comment}{\endcomment}
\theoremstyle{plain}
\newtheorem{theorem}[subsection]{Theorem}
\newtheorem{proposition}[subsection]{Proposition}
\newtheorem{lemma}[subsection]{Lemma}
\theoremstyle{definition}
\newtheorem{definition}[subsection]{Definition}
\newtheorem{example}[subsection]{Example}
\newtheorem{exercise}[subsection]{Exercise}
\newtheorem{situation}[subsection]{Situation}
\theoremstyle{remark}
\newtheorem{remark}[subsection]{Remark}
\newtheorem{remarks}[subsection]{Remarks}
\numberwithin{equation}{subsection}
\def\lim{\mathop{\mathrm{lim}}\nolimits}
\def\colim{\mathop{\mathrm{colim}}\nolimits}
\def\Spec{\mathop{\mathrm{Spec}}}
\def\Hom{\mathop{\mathrm{Hom}}\nolimits}
\def\Ext{\mathop{\mathrm{Ext}}\nolimits}
\def\SheafHom{\mathop{\mathcal{H}\!\mathit{om}}\nolimits}
\def\SheafExt{\mathop{\mathcal{E}\!\mathit{xt}}\nolimits}
\def\Sch{\mathit{Sch}}
\def\Mor{\mathop{\mathrm{Mor}}\nolimits}
\def\Ob{\mathop{\mathrm{Ob}}\nolimits}
\def\Sh{\mathop{\mathit{Sh}}\nolimits}
\def\NL{\mathop{N\!L}\nolimits}
\def\CH{\mathop{\mathrm{CH}}\nolimits}
\def\proetale{{pro\text{-}\acute{e}tale}}
\def\etale{{\acute{e}tale}}
\def\QCoh{\mathit{QCoh}}
\def\Ker{\mathop{\mathrm{Ker}}}
\def\Im{\mathop{\mathrm{Im}}}
\def\Coker{\mathop{\mathrm{Coker}}}
\def\Coim{\mathop{\mathrm{Coim}}}
\DeclareMathSymbol{\boxtimes}{\mathbin}{AMSa}{"02}
\def\QCohstack{\mathcal{QC}\!\mathit{oh}}
\def\Cohstack{\mathcal{C}\!\mathit{oh}}
\def\Spacesstack{\mathcal{S}\!\mathit{paces}}
\def\Quotfunctor{\mathrm{Quot}}
\def\Hilbfunctor{\mathrm{Hilb}}
\def\Curvesstack{\mathcal{C}\!\mathit{urves}}
\def\Polarizedstack{\mathcal{P}\!\mathit{olarized}}
\def\Complexesstack{\mathcal{C}\!\mathit{omplexes}}
\def\Pic{\mathop{\mathrm{Pic}}\nolimits}
\def\Picardstack{\mathcal{P}\!\mathit{ic}}
\def\Picardfunctor{\mathrm{Pic}}
\def\Deformationcategory{\mathcal{D}\!\mathit{ef}}
\usepackage{cleveref}
\crefname{originallemma}{original lemma}{original lemmas}
\Crefname{originallemma}{Original lemma}{Original lemmas}
\makeatletter
\newlabel{lemma-localize}{{08SF}{source}{}{sourcepointer-1}{}}
\newlabel{lemma-localize@cref}{{[originallemma][0][]08SF}{[1][0][]source}{}{}{}}
\makeatother
\crefname{originallemma}{original lemma}{original lemmas}
\Crefname{originallemma}{Original lemma}{Original lemmas}
\makeatletter
\newlabel{lemma-stalk-cotangent-complex}{{08T0}{source}{}{sourcepointer-2}{}}
\newlabel{lemma-stalk-cotangent-complex@cref}{{[originallemma][0][]08T0}{[1][0][]source}{}{}{}}
\makeatother
\crefname{originalremark}{original remark}{original remarks}
\Crefname{originalremark}{Original remark}{Original remarks}
\makeatletter
\newlabel{remark-map-sections-over-U}{{08SX}{source}{}{sourcepointer-3}{}}
\newlabel{remark-map-sections-over-U@cref}{{[originalremark][0][]08SX}{[1][0][]source}{}{}{}}
\makeatother
\crefname{originallemma}{original lemma}{original lemmas}
\Crefname{originallemma}{Original lemma}{Original lemmas}
\makeatletter
\newlabel{schemes-lemma-compare-constructions}{{01I7}{source}{}{sourcepointer-4}{}}
\newlabel{schemes-lemma-compare-constructions@cref}{{[originallemma][0][]01I7}{[1][0][]source}{}{}{}}
\makeatother
\crefname{originallemma}{original lemma}{original lemmas}
\Crefname{originallemma}{Original lemma}{Original lemmas}
\makeatletter
\newlabel{schemes-lemma-spec-sheaves}{{01HV}{source}{}{sourcepointer-5}{}}
\newlabel{schemes-lemma-spec-sheaves@cref}{{[originallemma][0][]01HV}{[1][0][]source}{}{}{}}
\makeatother
\crefname{originallemma}{original lemma}{original lemmas}
\Crefname{originallemma}{Original lemma}{Original lemmas}
\makeatletter
\newlabel{sheaves-lemma-stalk-pullback}{{008H}{source}{}{sourcepointer-6}{}}
\newlabel{sheaves-lemma-stalk-pullback@cref}{{[originallemma][0][]008H}{[1][0][]source}{}{}{}}
\makeatother
\title{The cotangent complex of a morphism of schemes restricted to an affine open is the sheafified ring-level complex}
\author{The Stacks Project Authors}
\date{stacks-project:cotangent.tex:local-snapshot-20261012; GNU FDL 1.2 or later}
\begin{document}
\maketitle
\setcounter{subsection}{1}
\noindent\textit{Editorial note.} All delivered mathematical text is the author's own, copied from the pinned Stacks Project chapter file cotangent.tex; no statement or proof step was rewritten, reordered or removed.\par\medskip
\noindent\textit{Editorial note.} The statement is about schemes rather than general algebraic spaces and includes the compatibility clause with smaller affine opens; both are inside the counted statement span.\par\medskip
\noindent\textit{Editorial note.} The six references of the counted spans are recorded as non-load-bearing navigation pointers; the argument that consumes them stays inside the counted proof span.\par\medskip
\noindent\textit{Editorial note.} Preamble: the retained preamble of this chapter (sources/stacks-cotangent/retained-preamble.tex, lines 1-63), containing the Xy-pic setup, the comment environments, the theorem declarations and the macros including \textbackslash{}NL.\par\medskip
\noindent\textit{Allowed standard prerequisites (editorial).}\begin{itemize}
\item The canonical map of complexes on sections over an open. Remark remark-map-sections-over-U, used at proof lines 4142-4147. Produces the first of the two maps.
\item The universal property of the sheafification functor. Schemes, Lemma schemes-lemma-compare-constructions, used at proof lines 4151-4153. Turns the module-level map into a map of sheaves.
\item The description of the stalks of the cotangent complex. Lemmas lemma-stalk-cotangent-complex and lemma-localize, used at proof lines 4154-4157. Show that the map is an isomorphism on stalks.
\item The description of the stalks of O\_X and f\textasciicircum{}\{-1\}O\_Y at a point. Schemes, Lemma schemes-lemma-spec-sheaves and Sheaves, Lemma sheaves-lemma-stalk-pullback, used at proof lines 4158-4164. Identify the two local rings appearing in the stalk computation.
\end{itemize}\medskip
\ifdefined\hypertarget\hypertarget{sourcepointer-1}{}\fi
\noindent\textit{Source reference.} Original lemma 08SF: focused-seeds/AGP/sources/stacks-snapshot-20261011/tags.txt, source lines 11128--11128. Navigation pointer into the Stacks Project; the standard fact it supplies is applied inside the counted proof span and the argument resting on it stays there.\par\medskip
\ifdefined\hypertarget\hypertarget{sourcepointer-2}{}\fi
\noindent\textit{Source reference.} Original lemma 08T0: focused-seeds/AGP/sources/stacks-snapshot-20261011/tags.txt, source lines 11148--11148. Navigation pointer into the Stacks Project; the standard fact it supplies is applied inside the counted proof span and the argument resting on it stays there.\par\medskip
\ifdefined\hypertarget\hypertarget{sourcepointer-3}{}\fi
\noindent\textit{Source reference.} Original remark 08SX: focused-seeds/AGP/sources/stacks-snapshot-20261011/tags.txt, source lines 11145--11145. Navigation pointer into the Stacks Project; the standard fact it supplies is applied inside the counted proof span and the argument resting on it stays there.\par\medskip
\ifdefined\hypertarget\hypertarget{sourcepointer-4}{}\fi
\noindent\textit{Source reference.} Original lemma 01I7: focused-seeds/AGP/sources/stacks-snapshot-20261011/tags.txt, source lines 1968--1968. Navigation pointer into the Stacks Project; the standard fact it supplies is applied inside the counted proof span and the argument resting on it stays there.\par\medskip
\ifdefined\hypertarget\hypertarget{sourcepointer-5}{}\fi
\noindent\textit{Source reference.} Original lemma 01HV: focused-seeds/AGP/sources/stacks-snapshot-20261011/tags.txt, source lines 1956--1956. Navigation pointer into the Stacks Project; the standard fact it supplies is applied inside the counted proof span and the argument resting on it stays there.\par\medskip
\ifdefined\hypertarget\hypertarget{sourcepointer-6}{}\fi
\noindent\textit{Source reference.} Original lemma 008H: focused-seeds/AGP/sources/stacks-snapshot-20261011/tags.txt, source lines 317--317. Navigation pointer into the Stacks Project; the standard fact it supplies is applied inside the counted proof span and the argument resting on it stays there.\par\medskip
\begin{lemma}
\label{lemma-morphism-affine-schemes}
Let $f : X \to Y$ be a morphism of schemes. Let $U = \Spec(B) \subset X$
and $V = \Spec(A) \subset Y$ be affine opens such that $f(U) \subset V$.
There is a canonical map
$$
\widetilde{L_{B/A}} \longrightarrow L_{X/Y}|_U
$$
of complexes which is an isomorphism in $D(\mathcal{O}_U)$.
This map is compatible with restricting to smaller affine opens
of $X$ and $Y$.
\end{lemma}

\begin{proof}
By Remark \ref{remark-map-sections-over-U}
there is a canonical map of complexes
$L_{\mathcal{O}_X(U)/f^{-1}\mathcal{O}_Y(U)} \to L_{X/Y}(U)$
of $B = \mathcal{O}_X(U)$-modules, which is compatible
with further restrictions. Using the canonical map
$A \to f^{-1}\mathcal{O}_Y(U)$ we obtain a canonical map
$L_{B/A} \to L_{\mathcal{O}_X(U)/f^{-1}\mathcal{O}_Y(U)}$
of complexes of $B$-modules.
Using the universal property of the $\widetilde{\ }$
functor (see Schemes, Lemma \ref{schemes-lemma-compare-constructions})
we obtain a map as in the statement of the lemma.
We may check this map is an isomorphism on cohomology sheaves
by checking it induces isomorphisms on stalks.
This follows immediately from
Lemmas \ref{lemma-stalk-cotangent-complex} and \ref{lemma-localize}
(and the description of the stalks of
$\mathcal{O}_X$ and $f^{-1}\mathcal{O}_Y$
at a point $\mathfrak p \in \Spec(B)$ as $B_\mathfrak p$ and
$A_\mathfrak q$ where $\mathfrak q = A \cap \mathfrak p$; references
used are Schemes, Lemma \ref{schemes-lemma-spec-sheaves}
and
Sheaves, Lemma \ref{sheaves-lemma-stalk-pullback}).
\end{proof}

\end{document}
